\subsection{Eingangsimpedanz bei AC-Kopplung}\label{sec:2-3}

\begin{highlightblock}[Berechnung der Eingangsimpedanz bei AC-Kopplung]
	Rechnerische Ermittlung des Betrages der gesamten Eingangsimpedanz (\autoref{fig:impedanz-oszi-ac-kopplung}) des verwendeten
	Oszilloskops mit angeschlossenem Kabel und AC–Kopplung für eine Kabelkapazität von
	$C_K = \SI{100}{\pico\farad}$ und eine Frequenz von $f = \SI{1}{\mega\hertz}$.

\end{highlightblock}

\subsubsection{Verwendetes Material}

\begin{table}[h]
	\centering
	\caption{Verwendetes Material}
	\begin{tabularx}{\textwidth}{XXX}
		\textbf{Messgerät} & \textbf{Verwendung} & \textbf{Anmerkungen}                                \\ \midrule
		TDS2014C           & Oszilloskop         & $Z_E = \SI{1}{\mega\ohm}\,||\,\SI{20}{\pico\farad}$ \\
		Agilent 33250A     & Funktionsgenerator  &                                                     \\
	\end{tabularx}
\end{table}

\subsubsection{Berechnung}

\begin{wrapfigure}[6]{r}{0.4\textwidth}
	\vspace{-32pt}\centering
	\begin{circuitikz}[scale=1.1]
		\ctikzset{bipoles/length=1cm}
		\draw (0,2) to[open, v=$U_E$, o-o] (0,0);
		\draw (0,2)
			-- ++(1,0) coordinate (k1)
			to[C, l=$C_{AC}$] ++(2,0) coordinate (k2)
			-- ++(1,0)
			to[R, l=$R_E$] ++(0,-2)
			-- (0,0); 
		\draw (k1) to[C, l=$C_K$, *-*] ++(0,-2);
		\draw (k2) to[C, l=$C_E$, *-*] ++(0,-2);
	\end{circuitikz}
	\caption{Eingangsimpedanz des Oszilloskopes bei AC-Kopplung}\label{fig:impedanz-oszi-ac-kopplung}
\end{wrapfigure}

Zur Berechnung der Gesamt-Eingangsimpedanz $Z$ betrachtet man das Ersatzbild in
\autoref{fig:impedanz-oszi-ac-kopplung}. $C_{AC}$ ist der errechnete Wert aus
\autoref{sec:2-3} \autoref{itm:2-P2-3}
\begin{itemize}[noitemsep]
	\item $R_E =\SI{1}{\mega\ohm}$
	\item $C_E =\SI{20}{\pico\farad}, C_K=\SI{100}{\pico\farad}, C_{AC}=\SI{22.7}{\nano\farad}$
	\item $\omega =2\pi\cdot \SI{1}{\mega\hertz}$
\end{itemize}

Es ergibt sich für die Gesamteingangsimpedanz $Z$ der Ausdruck (wobei $\jmath X = 1/\jmath\omega C$):
\begin{align}
	Z &= \jmath X_K \parallel \bigl(\jmath X_{AC} + (\jmath X_{E} \parallel R_E)\bigr) \nonumber\\
	&= \frac{
		1 + \jmath\omega R_E\bigl(C_{E} + C_{AC}\bigr)
	}{
		\omega^2 R_E\bigl(-C_E C_K - C_{AC} C_K - C_E C_{AC}\bigr) +
		\jmath \omega\bigl(C_K + C_{AC}\bigr)
	}
\label{eq:eingangswiderstand_oszi}
\end{align}

Ergebnis für eingesetzte Bauteilwerte:
\[
Z = (1.76-\jmath 1326.48)\si{\ohm} 
  = \SI{1326.48}{\ohm} \angle \SI{-89.92}{\degree}
  \implies |Z| \approx \SI{1346.5}{\ohm}
\]

\subsubsection{Erkenntnisse}

Die Eingangsimpedanz befindet sich hier im Kiloohm Bereich. Für empfindliche
Messungen könnte das bereits eine zu hohe belastung sein, und das Ergebnis
würde verfälscht werden.


